% !TeX root = ../main.tex
\section{研究方案}
\label{sec:research-plan}

\subsection{研究问题}
\label{sec:research-questions}

本研究以数字时代大学生人际交往方式的变化为切入点，将“大学生价值观”作为核心研究对象。研究重点考察“弱连接交往”这种情感投入低、人员流动快、信息来源杂、随时可替换的社交模式，如何改变大学生在学业职业、生活消费、人际交往、社会公共四个具体领域的价值取向，进而剖析这些实际变化与马克思主义历史发展方向的价值观、社会主义核心价值观以及当代中国政治与教育要求之间存在的矛盾与分歧，最终探索化解分歧、实现主流价值有效引领的可行办法。

随着算法推荐、匿名社区、临时群组以及“搭子”社交的普及，大学生的日常交往正从传统的熟人集体转向松散的弱连接网络。研究首先考察这种交往方式的变化在现实中让大学生的学业规划、日常消费、同伴互动以及公共参与发生了怎样的改变，这些具体观念与唯物史观的历史发展方向、社会主义核心价值观的要求存在哪些明显或隐蔽的差距。解释这些观念变化，需要考察媒介特征与心理活动如何参与其中。研究需要揭示弱连接交往在整体上如何借助庞杂信息的冲击、社交功能的拆解以及商业流量的诱导来动摇大学生原有的价值标准；同时在具体领域，深入分析它如何通过攀比心理、符号跟风、随时断联与情绪煽动等途径，分别作用于四个具体价值观维度，改变大学生的认知与选择。这些价值观变化也与大学生所处的社会环境和媒介条件有关。研究需要透过网络社交的表象，从我国社会转型期的生活状态变化、商业平台的流量盈利逻辑、大学生自身发展面临的现实压力，以及高校传统思政教育方式的不适应等多个角度，查明导致大学生具体价值观与主流价值要求产生背离的根本原因。面对弱连接交往日益普遍的现实，思政教育必须超越单纯的说教与被动限制。研究需要立足马克思主义的交往实践观与人的全面发展理论，探讨如何通过重建紧密的现实实践集体、改进主流价值的传播方式、提升大学生的马克思主义理论素养与信息辨别能力，引导青年将个人的日常追求融入民族复兴与国家现代化建设的大局之中。

判断弱连接交往中的价值分歧，需要建立“大学生具体价值观”与“主流应然价值观”的对比标准。以往研究大多笼统讨论网络对青年价值观的冲击，缺乏细致的分类参照。本研究将马克思主义关于生产力发展与劳动价值、消灭剥削与建立自由人联合体、人的本质是一切社会关系的总和、人民主体地位与人的自由全面发展等经典理论，同当代中国发展新质生产力、推进共同富裕、发展全过程人民民主、构建人类命运共同体以及落实立德树人“五育并举”的具体要求结合起来，转化为衡量大学生学业职业、生活消费、人际交往、社会公共四维价值观的明确标尺，理清价值分歧的边界。弱连接交往对主流价值认同的影响往往发生于日常社交之中，需要结合具体互动还原其作用过程。本研究结合社会网络理论与马克思主义政治经济学批判，理顺从媒介技术特征、社交行为模式到个体心理变化、价值观念偏移的完整链条，说明不同场景下弱连接影响机制的运作规律。教育引导还需要适应学生交往关系的变化。传统高校思政工作主要依靠院系、班级、宿舍等固定场所和熟人关系，面对分散、随时解散的弱连接社交往往力不从心。本研究尝试探索适应高流动网络社交环境的教育方法，让主流价值引领真正深入到学生的日常网络交往之中。

基于此，本研究的重点一方面在于搞清楚弱连接交往影响大学生四维具体价值观的作用过程与分歧表现，通过理论分析与实际调查的相互印证，细致拆解弱连接在四个具体领域引发观念偏差的机理；另一方面在于结合历史发展方向与国家人才培养要求，针对问题产生的原因，从集体建设、平台治理、理论教育等方面提出切实可行的应对策略。研究的难点则主要体现在两个方面：一是如何把抽象的价值哲学概念转化为可调查的具体指标。马克思主义价值观与社会主义核心价值观理论性很强，而大学生的日常弱连接社交体验非常细碎，在设计问卷和访谈提纲时，如何准确捕捉大学生在不知不觉中发生的观念变化，具有一定难度。二是如何区分弱连接交往与其他影响因素。大学生价值观的形成同时受到家庭背景、校园风气、就业形势等多种因素的影响，如何在分析中尽量排除其他因素的干扰，准确判断弱连接社交本身对四维价值观造成的实际影响，对研究设计提出了较高要求。

\subsection{研究内容}
\label{sec:research-content}

本研究以历史唯物主义与辩证唯物主义为指导，围绕弱连接交往的特征、影响机制、价值分歧及其成因展开研究，并据此探讨教育引导策略。

研究立足数字时代青年交往方式向弱连接转变的现实背景，通过梳理文献与理论思考，明确弱连接交往的当代表现，建立大学生四维具体价值观与主流价值观的对比框架。在这一框架下，先从理论层面提炼弱连接影响思想观念的三个共同机制；随后深入学业职业、生活消费、人际交往、社会公共四个具体领域，结合问卷调查、网络社群观察与访谈材料，按照“具体作用机制—实际观念变化—与主流价值分歧”的逻辑揭示矛盾的具体表现，并从社会环境、平台算法与学校教育三个角度分析分歧产生的深层原因。结合对作用机制与成因的分析，综合运用思想政治教育学、社会学与传播学知识，提出包含现实集体建设、网络内容优化与个人判断力提升在内的系统对策，促进青年个人发展与国家社会进步的统一。

研究首先通过梳理相关文献，考察弱连接交往的演变特征，明确主流价值的衡量标准，并界定相关概念：一方面从社会网络理论中的弱连接概念出发，说明当代大学生在算法平台、匿名社区、兴趣圈子及“搭子”社交中形成的低投入、脱离现实背景、高流动性的新型弱连接交往形态；另一方面建立价值观对比框架，界定以学业职业、生活消费、人际交往、社会公共为代表的“实际具体价值观”，以及以马克思主义历史发展方向价值观、社会主义核心价值观、当代中国政治与教育要求为代表的“主流应然价值观”。结合马克思关于人的社会关系本质、生产力决定作用、劳动发展史观以及人的自由全面发展等基本原理，分析交往方式与价值观念之间的关系。

弱连接交往对大学生认知与判断的影响，涉及信息来源、交往方式与平台推荐等因素。信息来源的扩展及其背景的缺失，是需要考察的作用过程。弱连接打破了熟人圈子的信息限制，让大量跨阶层、跨地域的信息零散地涌入学生视野；这种剥离了具体历史背景与真实劳动过程的信息展示，容易削弱大学生对社会发展规律的整体认识，动摇其对宏大理想的信服度。交往方式的变化则体现在社交功能的拆分与情感联系的减弱。弱连接将原本完整的同伴交往拆解为单一需求的临时匹配，如考研搭子、拼单群、游戏组队；随时可以退出的低成本断联方式降低了交往中的道德约束与责任感，助长了凡事算计得失的功利心态。商业算法推荐与视觉呈现也介入了这一过程。数字平台的商业算法深度介入弱连接社交，将享乐主义、消费攀比和极端个人主义包装成极具吸引力的画面与话题，在高频推送中潜移默化地改变学生的价值标准。

结合调查材料，研究进一步考察上述作用过程在学业职业、生活消费、人际交往、社会公共四个领域的表现，分析具体观念的变化及其与主流价值要求的分歧。

在学业职业价值观方面，重点分析“跨圈层盲目攀比与幸存者偏差放大机制”。弱连接网络频繁推送极端的暴富故事、优越生活展示，或是刻意渲染职场内卷与“躺平”言论，容易放大大学生的心理落差与对未来的恐慌。受此影响，大学生的学业职业观出现功利化与逃避风险的倾向，把专业学习单纯当成刷绩点和换文凭的工具，找工作时要么盲目追求短期赚快钱，要么消极退缩。这种变化背离了唯物史观重视劳动创造与生产力进步的发展方向，冲淡了核心价值观中“富强、敬业、诚信”的要求，也偏离了国家发展新质生产力、推进劳动教育以及号召青年投身战略产业和基层建设的育人导向。

在生活消费价值观方面，重点分析“视觉符号跟风与表面身份建构机制”。在缺乏深度了解的弱连接社交中，外在的消费标签成了大学生快速融入圈子、展示自我的捷径。在这种机制诱导下，大学生的消费观走向盲目追求品牌符号与即时享乐，热衷于网红消费、超前借贷与感官刺激，把个人价值寄托在买得起什么商品上。这使得青年陷入马克思所批判的商品拜物教与消费异化之中，对物质享受的过度追求挤压了精神成长的空间，偏离了马克思主义“人的自由全面发展”的目标，违背了核心价值观“文明、和谐”倡导的勤俭与环保理念，也与学校培养“德智体美劳”全面发展人才的要求背道而驰。

在人际交往价值观方面，重点分析“功能性切割与随时断联机制”。“搭子”式的弱连接严格以特定需求为界限，排斥深度的情感交流与长期的互助责任，把人际来往简化为精明的利益计算。长期处于这种模式下，大学生的人际交往观呈现出工具化与孤立化倾向，平时交往停留在表面客套与互利交换，遇到现实集体生活中的矛盾就习惯性地冷漠疏远或直接拉黑断联，退回封闭的个人世界。这种倾向割裂了“人的本质是一切社会关系的总和”的整体性，不利于培育面向“自由人联合体”的互助关系，削弱了当代中国推进共同富裕和构建人类命运共同体所需的集体主义精神，也使得核心价值观中的“平等、诚信、友善”只停留在表面礼貌，难以形成牢固的社会信任。

在社会公共价值观方面，重点分析“去历史化传播与情绪煽动机制”。弱连接网络的匿名性与高传染性，常常把复杂的社会公共事件剥离掉具体的历史脉络和制度背景，简化为非黑即白的对立标签与情绪发泄。受此裹挟，大学生的社会公共观容易滑向冷漠旁观与历史虚无倾向，表现为在网上跟风进行道德指责、对现实公共事务缺乏建设性参与，或是受片面信息误导产生偏激看法。这削弱了唯物史观所强调的人民群众创造历史的主动精神，背离了当代中国坚持以人民为中心、发展全过程人民民主的政治要求，对核心价值观中“民主、公正、法治、爱国”所要求的理性公民素养与制度自信造成了冲击。

上述价值分歧的形成，还需要联系社会环境、媒介生态与学校教育加以分析。在社会环境层面，我国市场经济深入发展带来的社会分工细化、人员高流动与个体独立意识增强，与社会主义集体主义、共同富裕的价值导向在青年步入社会的初期产生了一定摩擦，为功利化观念提供了现实土壤。在媒介生态层面，商业互联网平台受资本逐利驱动，算法推荐以争夺用户注意力和挑动情绪为首要目标，客观上纵容了炫富、贩卖焦虑和极端言论的传播，挤压了主流价值理性、系统传播的空间。在思政教育层面，高校传统的思政工作长期依赖行政班级、宿舍等熟人管理模式，宣传话语偏重宏观道理，对学生在网络社交中遇到的具体困惑回应不足，难以有效深入到流动、分散的弱连接社交圈中。

针对上述问题及其成因，研究从社会、高校与个人的联系出发，探讨集体建设、平台环境与理论教育中的引导措施。首先是重建紧密集体，用现实的“强连接共同体”稳住弱连接的过度流动，重点解决人际疏离与职业功利化问题。依托高校党团支部、科研团队、社会实践与志愿服务队等实体组织，在共同解决实际问题的集体劳动中，把浅层的“搭子”关系转化为有共同理想和深厚信任的同伴关系；同时利用弱连接信息面广的优势，搭建青年学生与国家重点行业、优秀建设者、基层工作者的交流渠道，用真实的国家建设成就引导青年将个人专业融入新质生产力发展。其次是优化平台生态与话语表达，重点解决消费跟风与公共情绪偏激问题。顺应网络传播特点，把马克思主义基本原理、国家发展成就和核心价值观转化为贴近学生生活的具体故事与优质文化内容；配合监管部门督促平台完善算法规则，抑制宣扬拜金享乐、制造择业焦虑与挑动群体对立的内容推荐。最后是加强理论方法教育，全面提升学生的价值判断力。在思政课教学中强化马克思主义政治经济学与唯物史观教育，帮助大学生看清商业资本逻辑、消费主义陷阱和算法推荐的运作本质，掌握透过碎片化信息看清社会发展规律的思维方法；完善校内民主参与和实践机制，引导大学生从网上的“情绪旁观者”转变为现实公共生活的“理性建设者”，在实践中培养担当意识。

\begin{thesisframework}
\begin{thesisoutline}
  \item \textbf{绪论}
  \begin{thesisoutline}
    \item 研究背景与研究意义
    \begin{thesisoutline}
      \item 研究背景
      \begin{thesisoutline}
        \item 大学生的交往环境发生了哪些变化？
        \item 这些联系为什么值得从价值观角度研究？
        \item 思想政治教育需要回应什么？
      \end{thesisoutline}
      \item 研究意义
      \begin{thesisoutline}
        \item 理论层面
        \item 实践层面
      \end{thesisoutline}
    \end{thesisoutline}
    \item 国内外研究综述
    \begin{thesisoutline}
      \item 国内研究综述
      \begin{thesisoutline}
        \item 关于“弱连接”的研究
        \begin{thesisoutline}
          \item 弱连接的概念界定与关系特征
          \item 弱连接的社会功能及其作用条件
          \item 大学生弱连接交往的表现形式
        \end{thesisoutline}
        \item 关于大学生价值观的研究
        \item 关于“弱连接”与价值观的关系研究
      \end{thesisoutline}
      \item 国外研究综述
      \begin{thesisoutline}
        \item 关于“弱连接”（weak ties）的研究
        \item 关于青年价值观与公民教育的研究
      \end{thesisoutline}
      \item 研究述评
    \end{thesisoutline}
    \item 研究思路与方法
    \begin{thesisoutline}
      \item 研究思路
      \item 研究方法
      \par 访谈法（半结构化访谈）
    \end{thesisoutline}
    \item 研究重难点与创新点
    \begin{thesisoutline}
      \item 重难点
      \item 创新点
    \end{thesisoutline}
  \end{thesisoutline}
  \item \textbf{弱连接与大学生价值观的理论阐释}
  \begin{thesisoutline}
    \item 核心概念阐释
    \begin{thesisoutline}
      \item 弱连接的内涵及其特征
      \begin{thesisoutline}
        \item 弱连接的基本内涵
        \item 弱连接与强连接的区别及联系
        \item 大学生弱连接交往的主要特征
      \end{thesisoutline}
      \item 大学生价值观的内涵
      \begin{thesisoutline}
        \item 价值观的基本内涵
        \item 大学生价值观形成和确立的基本规律
        \item 大学生价值观的影响因素
      \end{thesisoutline}
    \end{thesisoutline}
    \item 理论基础
    \begin{thesisoutline}
      \item 马克思主义价值观理论
      \item 马克思主义人的全面发展理论
      \item 格兰诺维特弱关系理论
    \end{thesisoutline}
  \end{thesisoutline}
  \item \textbf{大学生弱连接问题的基本状况}
  \begin{thesisoutline}
    \item 调查设计
    \begin{thesisoutline}
      \item 调查目的与主要内容
      \item 研究方法的选择依据
      \item 调查对象的选择与弱连接的识别
      \item 访谈提纲与问题设置
    \end{thesisoutline}
    \item 调查实施与材料分析
    \begin{thesisoutline}
      \item 访谈开展与材料收集
      \item 访谈材料的整理与分类
      \item 访谈材料的比较与解释
    \end{thesisoutline}
    \item 受访大学生弱连接交往的主要表现
    \begin{thesisoutline}
      \item 交往对象与关系建立途径
      \item 交往动机与主要交流内容
      \item 交往投入与关系维系方式
      \item 关系变化与交往体验
    \end{thesisoutline}
    \item 受访大学生弱连接交往中的价值表达
    \begin{thesisoutline}
      \item 人生目标与成功评价的表达
      \item 诚信公平与友善互助的判断
      \item 职业意义与择业标准的认识
    \end{thesisoutline}
    \item 大学生弱连接交往中的价值议题
    \begin{thesisoutline}
      \item 关于人生目标与成功评价的议题
      \item 关于个人选择与集体生活的议题
      \item 关于交往意义与相互帮助的议题
    \end{thesisoutline}
  \end{thesisoutline}
  \item \textbf{弱连接对大学生价值观的影响}
  \begin{thesisoutline}
    \item 大学生价值观的内容领域与本文考察重点
    \begin{thesisoutline}
      \item 大学生价值观的主要内容领域
      \item 弱连接影响价值观的考察范围
      \item 考察领域的识别依据与边界
    \end{thesisoutline}
  \end{thesisoutline}
  \item \textbf{弱连接影响大学生价值观的机制与条件}
  \begin{thesisoutline}
    \item 弱连接影响大学生人生价值观的机制与条件
    \begin{thesisoutline}
      \item 生活经验接触：影响人生价值判断的认识基础
      \item 社会比较参照：影响人生目标与成功标准的评价依据
      \item 自我发展预期：影响人生价值追求的现实判断
    \end{thesisoutline}
    \item 弱连接影响大学生道德价值观的机制与条件
    \begin{thesisoutline}
      \item 交往规范感知：影响道德判断的情境依据
      \item 互动结果归因：影响对他人行为的道德评价
      \item 道德情感体验：影响道德要求的接受与认同
    \end{thesisoutline}
    \item 弱连接影响大学生职业价值观的机制与条件
    \begin{thesisoutline}
      \item 职业信息认识：影响择业评价的信息依据
      \item 职业能力判断：影响职业选择的可行性认识
      \item 职业结果预期：影响职业价值的权衡与排序
    \end{thesisoutline}
  \end{thesisoutline}
  \item \textbf{弱连接交往情境下大学生价值观教育的应对策略}
  \begin{thesisoutline}
    \item 深化人生价值观教育，引导大学生确立合理的价值追求
    \begin{thesisoutline}
      \item 拓展人生经验参照，增进对多样发展道路的理解
      \item 加强人生价值辨析，克服单一成功标准的局限
      \item 联系个人成长实际，促进个人价值与社会价值相统一
    \end{thesisoutline}
    \item 深化道德价值观教育，引导大学生形成正确的交往判断
    \begin{thesisoutline}
      \item 结合真实交往情境，深化诚信与公平观念教育
      \item 加强交往经历反思，提高道德判断与辨析能力
      \item 丰富合作互助体验，促进道德认知与行为践行相结合
    \end{thesisoutline}
    \item 深化职业价值观教育，引导大学生形成合理的择业标准
    \begin{thesisoutline}
      \item 整合多元职业信息，提升职业认识的真实性与全面性
      \item 加强职业价值辨析，引导择业标准的合理权衡
      \item 结合职业探索指导，促进个人志向与社会需要衔接
    \end{thesisoutline}
  \end{thesisoutline}
\end{thesisoutline}
\par\smallskip\noindent\textbf{结语}
\end{thesisframework}


\subsection{研究方法}
\label{sec:research-methods}

本研究坚持马克思主义辩证唯物主义与历史唯物主义的根本方法。一方面坚持“社会存在决定社会意识”的原则，把大学生价值观的具体变化放在数字技术发展、平台商业运作与社会交往方式改变的现实条件下去考察，避免陷入单纯道德说教的误区。另一方面坚持辩证分析方法，既看到弱连接交往在拓宽学生视野、提供便利方面的积极作用，也正视其在冲淡深层价值认同方面的消极影响，以理论分析为主导，辅以适度的调查与案例分析，确保研究结论扎实可靠。

研究以文献与理论分析为主，结合网络观察、问卷调查与代表性访谈。

研究将重点精读《马克思恩格斯文集》中的《1844年经济学哲学手稿》《德意志意识形态》《共产党宣言》等经典著作，梳理其中关于异化劳动、商品拜物教、社会交往形式、人的本质以及自由人联合体的核心论述；同时认真研读党的创新理论文献以及国家关于发展新质生产力、推进共同富裕、全过程人民民主、社会主义核心价值观建设与高校思政课建设的重要文件，以此确立衡量价值观的主流标准。在此基础上，广泛检索国内外学术数据库中关于弱连接理论、青年网络社交、“搭子”文化及大学生价值观变化的学术论文，理清现有研究的成果与不足，运用马克思主义立场观点方法对这些跨学科材料进行理论整合与逻辑推演，构建起本文的分析框架。

为使理论分析具备现实支撑，研究将采用网络观察与文本分析法。针对弱连接交往高度依赖特定社交平台的特点，选取大学生常用来找“搭子”、分享消费、讨论学业与公共话题的代表性平台（如小红书、豆瓣小组、知乎及校园论坛）进行持续观察。研究者通过浏览记录这些社区里的常见话题、互动习惯以及关系建立和解散的过程，收集涉及职业焦虑、消费跟风、人际疏离等典型讨论帖与高赞评论。通过对这些具体案例和文本内容的归类整理，直观展现跨圈层攀比、符号消费、随时断联与情绪宣泄在网络社交中的实际表现，作为理论分析的现实佐证。

在此基础上，通过小规模问卷调查与代表性访谈了解大学生的交往经历和具体认识。考虑到个人研究精力与可操作性，问卷调查定位于摸底和辅助验证。围绕大学生的弱连接社交习惯（如平台使用频率、找“搭子”经历、陌生人信息接触情况）以及学业职业、生活消费、人际交往、社会公共四个维度的具体观念设计调查问卷，利用线上问卷工具在本校及周边高校发放并回收200份左右有效问卷，通过基础的数据统计来反映大学生弱连接交往与价值观变化的总体趋势。同时，在身边同学及思政课教师、辅导员中选取10至15人开展半结构化访谈，重点了解学生在弱连接社交中的真实感受、面对网络庞杂信息时的心理波动以及对主流价值要求的看法。访谈记录与问卷数据、网络典型案例相互补充，共同支撑起前文的机制分析与对策建议。

\subsection{创新之处}
\label{sec:innovation}

本研究以唯物史观关于人的社会关系的论述为依据，从青年人际交往关系的变化出发，考察弱连接交往如何影响大学生的价值判断。研究将大学生价值观细化为学业职业、生活消费、人际交往、社会公共四个日常领域，结合马克思主义相关理论、社会主义核心价值观及当代中国政治与教育要求建立分析参照，尝试梳理不同领域中交往方式、具体观念与价值分歧之间的联系，并分析相应的作用机制。

围绕上述问题，研究将马克思主义经典理论分析与网络社交案例观察、问卷调查和访谈相结合，使理论讨论能够得到具体交往材料的支持，也使对策建议与大学生的实际生活及所发现的问题相对应。这一研究安排旨在对弱连接交往的价值影响作出更具体的解释，为高校开展相关教育引导提供参考。
